\chapter{Sprint 1 : Accès et gestion du compte utilisateur}
\label{chap:sprint1}

\section{Introduction}
\addcontentsline{toc}{section}{Introduction}

Avant de pouvoir faire quoi que ce soit d’intelligent, il faut d’abord savoir à qui on parle. C’est le rôle du premier sprint : mettre en place le module d’authentification et la gestion du compte utilisateur.

Ce n’est pas la partie la plus spectaculaire du projet, mais c’est une fondation indispensable. Toutes les données collectées par le système : sessions de travail, scores de concentration, planning, documents, sont liées à un compte utilisateur. Sans authentification fiable, rien de tout ça ne tient. Ce sprint a également été l’occasion de mettre en place les premières interfaces de l’application mobile et d’établir la communication entre le frontend Flutter et le backend FastAPI.

\section{Backlog du Sprint 1}

Le backlog du Sprint 1 couvre les quatre fonctionnalités de base liées à la gestion du compte utilisateur. Chaque user story est accompagnée de la tâche technique correspondante, telle qu’elle a été planifiée au début du sprint.

Le tableau~\ref{tab:backlogs1} présente le détail de ce backlog.

\begin{table}[H]
\centering
\renewcommand{\arraystretch}{1.3}
\begin{tabular}{|c|p{3cm}|p{6cm}|p{4cm}|}
\hline
\textbf{ID} & \textbf{Fonctionnalité} & \textbf{User Story} & \textbf{Tâche} \\
\hline

US1 & S'inscrire 
& En tant qu’utilisateur, je souhaite créer un compte afin d’accéder aux fonctionnalités de l’application 
& Développement des interfaces mobiles et implémentation de l’API d’inscription \\ \hline

US2 & Se connecter
& En tant qu’utilisateur, je souhaite me connecter de manière sécurisée à mon compte
& Mise en place d’un mécanisme de connexion sécurisée avec gestion des sessions utilisateur \\ \hline

US3 & Gérer son profil
& En tant qu’utilisateur, je souhaite consulter et modifier mes informations personnelles
& Développement des fonctionnalités de consultation et de modification du profil utilisateur \\ \hline

US4 & Se déconnecter
& En tant qu’utilisateur, je souhaite me déconnecter afin de sécuriser ma session
& Gestion de la déconnexion et fermeture de la session utilisateur \\ \hline

\end{tabular}
\caption{Backlog du Sprint 1}
\label{tab:backlogs1}
\end{table}



\section{Spécifications fonctionnelles}

\subsection{Diagramme de cas d’utilisation du Sprint 1}
Le diagramme suivant représente les interactions entre l’utilisateur et le système pour les fonctionnalités développées durant ce sprint.
\begin{figure}[H]
\centering
\includegraphics[width=1\textwidth]{images/Sprint1UC.png}
\caption{Diagramme de cas d'utilisation du Sprint 1}
\end{figure}
%%%%%%%%%%%%%%%%%%%%%%%%%%%%%%%%%%%%%%%%%%%%%%%%%%%%%%%%%%%

\subsection{Description textuelle des cas d’utilisation}

Chaque cas d’utilisation est décrit sous forme de fiche synthétique qui détaille les conditions de déclenchement, les étapes du scénario nominal et les situations d’erreur possibles.

%----------------------------------------
\subsubsection{Cas d’utilisation : S’inscrire}

L’inscription est la première action que tout nouvel utilisateur doit effectuer. Elle implique la saisie de ses informations, leur validation côté serveur, et la création d’un compte sécurisé. Le tableau~\ref{tab:uc_register} en décrit les étapes.

\begin{table}[H]
\centering
\renewcommand{\arraystretch}{1.2}
\begin{tabular}{|p{4cm}|p{10cm}|}
\hline
\textbf{Acteur} & Utilisateur \\ \hline

\textbf{Objectif} & Créer un compte pour accéder aux fonctionnalités du système \\ \hline

\textbf{Conditions initiales} & Aucun compte existant avec le même email \\ \hline

\textbf{Déroulement} &
\begin{itemize}
    \item L’utilisateur ouvre l’écran d’inscription
    \item L’utilisateur saisit ses informations personnelles
    \item Le système vérifie les données saisies
    \item Le système crée le compte utilisateur
    \item Une session est automatiquement ouverte
\end{itemize} \\ \hline

\textbf{Cas d’erreur} & Adresse email déjà utilisée → refus de création du compte \\ \hline

\textbf{Résultat} & Compte créé et accès direct à l’application \\ \hline
\end{tabular}
\caption{Description du cas d’utilisation « S’inscrire »}
\label{tab:uc_register}
\end{table}

%----------------------------------------
\subsubsection{Cas d’utilisation : Se connecter}

La connexion permet à un utilisateur déjà inscrit d’accéder à son compte. Le Tableau~\ref{tab:uc_login} décrit ce processus.

\begin{table}[H]
\centering
\renewcommand{\arraystretch}{1.2}
\begin{tabular}{|p{4cm}|p{10cm}|}
\hline
\textbf{Acteur} & Utilisateur \\ \hline

\textbf{Objectif} & Accéder à son compte de manière sécurisée \\ \hline

\textbf{Conditions initiales} & Compte utilisateur valide \\ \hline

\textbf{Déroulement} &
\begin{itemize}
    \item L’utilisateur saisit son adresse email et son mot de passe
    \item Le système vérifie la correspondance des identifiants
    \item Le système ouvre une session sécurisée
    \item L’utilisateur accède à l’application
\end{itemize} \\ \hline

\textbf{Cas d’erreur} & Identifiants incorrects → accès refusé \\ \hline

\textbf{Résultat} & Session utilisateur active \\ \hline
\end{tabular}
\caption{Description du cas d’utilisation « Se connecter »}
\label{tab:uc_login}
\end{table}

%----------------------------------------
\subsubsection{Cas d’utilisation : Se connecter via Gmail}

Pour simplifier l’accès, le système offre une option de connexion via Google. L’utilisateur n’a pas besoin de créer un mot de passe ; il s’authentifie directement avec son compte Gmail. Le tableau~\ref{tab:uc_login_google} décrit le déroulement de ce processus.

\begin{table}[H]
\centering
\renewcommand{\arraystretch}{1.2}
\begin{tabular}{|p{4cm}|p{10cm}|}
\hline
\textbf{Acteur} & Utilisateur \\ \hline

\textbf{Objectif} & Se connecter au système via un compte Google \\ \hline

\textbf{Conditions initiales} & Compte Google valide avec une adresse email vérifiée \\ \hline

\textbf{Déroulement} &
\begin{itemize}
    \item L’utilisateur sélectionne l’option « Se connecter avec Google »
    \item L’utilisateur est redirigé vers la page de connexion Google
    \item L’utilisateur autorise l’application à accéder à son compte Google
    \item Le système transmet les informations d’autorisation pour vérification
    \item Le système vérifie l’autorisation auprès des serveurs Google
    \item Le système vérifie que l’adresse email est valide et confirmée
    \item Le système recherche le compte ; il le crée automatiquement si l’utilisateur n’est pas encore inscrit
    \item Une session est ouverte et l’utilisateur accède à l’application
\end{itemize} \\ \hline

\textbf{Cas d’erreur} & Autorisation Google invalide ou expirée → accès refusé \newline Adresse email non confirmée par Google → accès interdit \\ \hline

\textbf{Résultat} & Session utilisateur active \\ \hline
\end{tabular}
\caption{Description du cas d’utilisation « Se connecter via Gmail »}
\label{tab:uc_login_google}
\end{table}

%----------------------------------------
\subsubsection{Cas d’utilisation : Gérer son profil}

Une fois connecté, l’utilisateur peut consulter et modifier ses informations personnelles ainsi que ses préférences. Le tableau~\ref{tab:uc_profile} décrit ce cas d’utilisation.

\begin{table}[H]
\centering
\renewcommand{\arraystretch}{1.2}
\begin{tabular}{|p{4cm}|p{10cm}|}
\hline
\textbf{Acteur} & Utilisateur \\ \hline

\textbf{Objectif} & Consulter et modifier ses informations personnelles \\ \hline

\textbf{Conditions initiales} & Utilisateur authentifié \\ \hline

\textbf{Déroulement} &
\begin{itemize}
    \item L’utilisateur ouvre sa page de profil
    \item Le système affiche les informations enregistrées
    \item L’utilisateur modifie les informations souhaitées
    \item Le système valide et enregistre les changements
\end{itemize} \\ \hline

\textbf{Cas d’erreur} & Données non valide → refus de modification des données \\ \hline

\textbf{Résultat} & Données mises à jour \\ \hline
\end{tabular}
\caption{Description du cas d’utilisation « Gérer son profil »}
\label{tab:uc_profile}
\end{table}

%----------------------------------------
\subsubsection{Cas d’utilisation : Se déconnecter}

La déconnexion met fin à la session en cours et protège l’accès en supprimant les données de connexion locales. Le tableau~\ref{tab:uc_logout} en résume les étapes.

\begin{table}[H]
\centering
\renewcommand{\arraystretch}{1.2}
\begin{tabular}{|p{4cm}|p{10cm}|}
\hline
\textbf{Acteur} & Utilisateur \\ \hline

\textbf{Objectif} & Mettre fin à la session utilisateur \\ \hline

\textbf{Conditions initiales} & Utilisateur connecté \\ \hline

\textbf{Déroulement} &
\begin{itemize}
    \item L’utilisateur sélectionne l’option de déconnexion
    \item Le système ferme la session en cours
    \item Les données de connexion locales sont supprimées
    \item L’utilisateur est redirigé vers l’écran de connexion
\end{itemize} \\ \hline

\textbf{Résultat} & Session terminée et accès sécurisé \\ \hline
\end{tabular}
\caption{Description du cas d’utilisation « Se déconnecter »}
\label{tab:uc_logout}
\end{table}

%%%%%%%%%%%%%%%%%%%%%%%%%%%%%%%%%%%%%%%%%%%%%%%%%%%%%%%%%%%%%%%%%%%%%%
\section{Conception}
\subsection{Diagrammes de séquence}

Les diagrammes de séquence ci-dessous décrivent les échanges entre les composants du système pour les principaux cas d’utilisation du sprint.

%----------------------------------------
\subsubsection{Processus d’inscription}

Lors de l’inscription, l’application transmet les données saisies au serveur, qui vérifie d’abord si l’adresse email n’est pas déjà utilisée. Si tout est en ordre, le mot de passe est haché avant d’être stocké, le compte est créé, et une session est ouverte automatiquement. L’utilisateur est directement connecté sans avoir à se reconnecter.

\begin{figure}[H]
\centering
\includegraphics[width=0.85\textwidth]{images/inscriptionSeq.png}
\caption{Diagramme de séquence du processus d’inscription}
\label{fig:seq_register}
\end{figure}

%----------------------------------------
\subsubsection{Processus de connexion}

Pour la connexion, l’application envoie les identifiants au backend, qui compare le mot de passe saisi à son empreinte stockée via bcrypt. Si la vérification réussit, deux tokens sont générés : un access token de courte durée et un refresh token, et la session est ouverte.

\begin{figure}[H]
\centering
\includegraphics[width=0.85\textwidth]{images/AuthViaEmail.png}
\caption{Diagramme de séquence du processus d’authentification}
\label{fig:seq_login}
\end{figure}

%----------------------------------------


%----------------------------------------
\subsubsection{Gestion du profil utilisateur}

L’accès au profil nécessite que l’utilisateur soit authentifié. Le système récupère les données depuis la base, les affiche, et enregistre les éventuelles modifications après validation.

\begin{figure}[H]
\centering
\includegraphics[width=0.85\textwidth]{images/GérerProfilSeq.png}
\caption{Diagramme de séquence de la gestion du profil}
\label{fig:seq_profile}
\end{figure}

%----------------------------------------
%\subsubsection{Processus de déconnexion}

%Le processus de déconnexion permet de mettre fin à la session active de l’utilisateur afin de garantir la sécurité de l’accès au système. Comme illustré dans la Figure~\ref{fig:seq_logout}, l’application envoie une requête de déconnexion au backend accompagnée du token d’authentification.Le système invalide ce token, empêchant ainsi toute utilisation ultérieure, puis confirme la déconnexion. L’utilisateur est ensuite redirigé vers l’écran de connexion, assurant ainsi une terminaison propre et sécurisée de la session.

%\begin{figure}[H]
%\centering
%\includegraphics[width=0.85\textwidth]{images/logout_sequence.png}
%\caption{Diagramme de séquence du processus de déconnexion}
%\label{fig:seq_logout}
%\end{figure}

\subsection{Diagramme de classe du sprint 1 }
Le diagramme de classes de la Figure~\ref{fig:class_sprint1} montre comment le module d’authentification est structuré. La classe \textit{User} regroupe les informations de compte, tandis que \textit{UserProfile} stocke les préférences personnalisées de l’utilisateur. \textit{AuthToken} gère les tokens de session, et \textit{AuthService} centralise la logique d’inscription et de connexion. Cette séparation des responsabilités simplifie la maintenance et renforce la sécurité du module.

\begin{figure}[H]
\centering
\includegraphics[width=0.75\textwidth]{images/DiagClassSprint1.png}
\caption{Diagramme de classes du module d’authentification et de gestion utilisateur}
\label{fig:class_sprint1}
\end{figure}
%%%%%%%%%%%%%%%%%%%%%%%%%%%%%%%%%%%%%%%%%%%%%%%%%%%%%%%%%%%%%%%%
%----------------------------------------
\section{Environnement technique}

Cette section détaille les technologies utilisées dans le Sprint 1. Les choix ont été orientés par des impératifs de sécurité, notamment pour la gestion des mots de passe et des sessions, ainsi que par la nécessité d'une architecture maintenable sur le long terme.

%----------------------------------------
\subsection{JSON Web Token (JWT)}

On avait d'abord envisagé un système de sessions côté serveur avec Redis ; c'est une approche courante, et on la connaissait bien. Mais ajouter Redis aurait alourdi l'infrastructure au moment où on essayait justement de la simplifier. JWT s'est imposé comme un meilleur compromis : pas d'état à maintenir côté serveur, une validation possible sans appel base de données, et une gestion propre de l'expiration.

Les JWT \cite{jwt} (RFC 7519) gèrent les sessions utilisateur. À la connexion, le serveur génère deux tokens :
\begin{itemize}
    \item \textbf{Access token} : valide 15 minutes, il accompagne chaque requête de l'application pour authentifier l'utilisateur.
    \item \textbf{Refresh token} : valide 7 jours, il permet de renouveler l'access token expirée sans redemander les identifiants.
\end{itemize}

Chaque token est structuré en trois parties encodées en Base64 : l'en-tête, la charge utile (contenant les données utilisateur), et la signature garantissant l'intégrité du token.

%----------------------------------------
\subsection{Hachage des mots de passe avec Bcrypt}

On a regardé argon2id brièvement ; c'est l'algorithme recommandé par OWASP depuis quelques années. Mais l'intégration Python était un peu plus complexe, et bcrypt \cite{bcrypt} s'intégrait mieux avec \texttt{passlib}, qu'on utilisait déjà. Pour un projet de cette taille, le gain de sécurité marginal d'argon2 ne justifiait pas la complexité supplémentaire.

bcrypt est conçu pour être lent, ce qui le rend résistant aux attaques par force brute. Il intègre aussi un mécanisme de salage : même deux utilisateurs avec le même mot de passe auront des empreintes différentes en base de données. Le facteur de coût a été fixé à 12 après quelques tests : à 10, ça allait trop vite ; à 14, la connexion prenait un temps perceptible sur le serveur de dev. 12 était le bon milieu. Lors de l'authentification, bcrypt compare directement le mot de passe saisi à son empreinte, sans jamais le décoder.

%----------------------------------------
\subsection{Authentification via Google OAuth 2.0}

L'ajout de Google OAuth n'était pas prévu au départ ; on pensait que l'authentification classique email/mot de passe suffirait. Mais lors des premiers tests avec des utilisateurs réels, la question "est-ce qu'on peut se connecter avec Google ?" est revenue plusieurs fois. Les étudiants utilisent déjà leurs comptes Google pour tout le reste : Drive, Meet, mails académiques. Ajouter une connexion supplémentaire créait une friction inutile. On a donc décidé d'intégrer OAuth 2.0 dès le Sprint 1.

Google OAuth 2.0 permet à l'utilisateur de s'authentifier avec son compte Google, sans créer de nouveau mot de passe. Le flux se déroule en trois étapes :
\begin{enumerate}
    \item L'application mobile redirige l'utilisateur vers la page de connexion Google, qui lui demande d'autoriser l'accès.
    \item Google retourne un \textit{id\_token} signé contenant les informations du compte (email, nom, photo de profil).
    \item Le serveur vérifie l'authenticité du token auprès des serveurs Google, extrait l'adresse email, et crée le compte utilisateur si celui-ci n'existe pas encore.
\end{enumerate}

%----------------------------------------
\subsection{FastAPI}

FastAPI est un framework web Python moderne conçu pour la construction d'APIs RESTful avec de hautes performances. Il repose sur le standard ASGI et intègre nativement la validation des données via Pydantic ainsi que la génération automatique de documentation interactive (Swagger UI).

Dans ce sprint, FastAPI expose les routes d'authentification : inscription, connexion, connexion via Google, gestion du profil et renouvellement de session.

%----------------------------------------
\subsection{Flutter et Dio}

Flutter est le framework de développement mobile utilisé pour construire l'application cliente. La communication avec le serveur est assurée par la bibliothèque \textit{Dio}, un client HTTP pour Dart qui permet de gérer les requêtes HTTP, les intercepteurs et la gestion automatique des erreurs réseau.

Un intercepteur JWT est configuré dans Dio pour ajouter automatiquement l'access token à chaque requête, et pour déclencher le renouvellement de session lorsque ce token expire.

%----------------------------------------
\subsection{SQLAlchemy et PostgreSQL}

Les données utilisateur sont persistées dans une base de données PostgreSQL via l'ORM SQLAlchemy. Deux tables principales sont utilisées dans ce sprint :
\begin{itemize}
    \item \textbf{users} : stocke les informations de compte (email, empreinte du mot de passe, rôle, statut).
    \item \textbf{user\_profiles} : stocke les préférences personnalisées (objectif de focus quotidien, paramètres de notification).
\end{itemize}

SQLAlchemy permet de définir les modèles de données sous forme de classes Python et de gérer les relations entre tables de manière transparente.

%----------------------------------------
\section{Implémentation}

Cette section présente les interfaces réalisées durant le Sprint 1, couvrant l’onboarding, l’authentification et la gestion du profil utilisateur.

%--------------------------------------

\subsection{Interfaces de l’application mobile}

Les interfaces du Sprint 1 couvrent trois niveaux : l’accueil de l’application, les écrans d’authentification, et la gestion du profil.

%----------------------------------------
\subsubsection{Écrans d’accueil et d’introduction}

Les écrans d’accueil servent d’introduction au système. Ils présentent brièvement le concept de Smart Focus et orientent l’utilisateur vers l’inscription ou la connexion.

\begin{figure}[H]
\centering
\includegraphics[width=0.35\textwidth]{images/1.jpg}
\hspace{0.5cm}
\includegraphics[width=0.35\textwidth]{images/2.jpg}
\caption{Écrans d’accueil et d’introduction}
\label{fig:welcome_onboarding}
\end{figure}

%----------------------------------------
\subsubsection{Interfaces d’authentification}

Ces écrans gèrent la connexion et l’inscription. La connexion via Google a également été intégrée pour simplifier l’accès.

\begin{figure}[H]
\centering
\includegraphics[width=0.35\textwidth]{images/3.jpg}
\hspace{0.5cm}
\includegraphics[width=0.35\textwidth]{images/4.jpg}
\caption{Interfaces de connexion et d’inscription}
\label{fig:auth}
\end{figure}

%----------------------------------------
\subsubsection{Interface de profil et paramètres}

L’écran de profil permet à l’utilisateur de consulter et modifier ses informations personnelles, et de configurer ses préférences : objectif de focus quotidien, créneaux horaires préférés pour les sessions, notifications.

\begin{figure}[H]
\centering
\includegraphics[width=0.35\textwidth]{images/5.jpg}
\hspace{0.5cm}
\includegraphics[width=0.35\textwidth]{images/6.jpg}
\caption{Interface de profil et paramètres}
\label{fig:profile_settings}
\end{figure}

%----------------------------------------
\subsubsection{Interface du tableau de bord}

Le tableau de bord constitue la page principale de l’application. Il offre à l’utilisateur une vue synthétique de son état courant : session de travail active, score de concentration en temps réel, progression du planning du jour et résumé des dernières activités.

\begin{figure}[H]
\centering
% TODO: remplacer dashboard1.jpg et dashboard2.jpg par les vraies captures
\includegraphics[width=0.35\textwidth]{images/dashboard1.jpg}
\hspace{0.5cm}
\includegraphics[width=0.35\textwidth]{images/dashboard2.jpg}
\caption{Interface du tableau de bord}
\label{fig:dashboard}
\end{figure}

%----------------------------------------
\subsubsection{Interface des statistiques}

L’écran des statistiques présente les indicateurs de performance agrégés sur la période sélectionnée : temps d’étude total, taux de complétion des sessions, évolution du score de concentration et sujets les plus travaillés. Ces données permettent à l’utilisateur d’identifier ses points forts et ses axes d’amélioration.

\begin{figure}[H]
\centering
% TODO: remplacer stats1.jpg et stats2.jpg par les vraies captures
\includegraphics[width=0.35\textwidth]{images/stats1.jpg}
\hspace{0.5cm}
\includegraphics[width=0.35\textwidth]{images/stats2.jpg}
\caption{Interface des statistiques de performance}
\label{fig:statistics}
\end{figure}

%----------------------------------------
\subsubsection{Interface de suivi du sommeil}

L’écran de sommeil permet à l’utilisateur d’enregistrer ses nuits et de consulter l’historique de ses données de sommeil : durée totale, score de qualité et corrélation avec les performances d’étude. Ces informations alimentent directement le planificateur intelligent pour adapter l’intensité des sessions du lendemain.

\begin{figure}[H]
\centering
% TODO: remplacer sleep1.jpg et sleep2.jpg par les vraies captures
\includegraphics[width=0.35\textwidth]{images/sleep1.jpg}
\hspace{0.5cm}
\includegraphics[width=0.35\textwidth]{images/sleep2.jpg}
\caption{Interface de suivi du sommeil}
\label{fig:sleep_tracking}
\end{figure}

%%%%%%%%%%%%%%%%%%%%%%%%%%%%%%%%%%%%%%%%%%%%%%%%%%%%%%%%%%%%%%%
%%%%%%%%%%%%%%%%%%%%%%%%%%%%%%%%%%%%%%%%%%%%%%%%%%%%%%%%%%%%%%%%%%%%%%%%%%%%%%%%%%%%%%%%%%%%%%%%%%%%%%%%%%%%%%%%%%%%%%%%%%%%%%%%%%%%%%%%%%%%%%%%%%%%%%%%%%%%%%%%%%%%%%%%%%%%%%%%%%%%%%%%%%%%%%%%%%%%%%%%%%%%%%%%%%%%%%%%%%%%%%%%%%%%%%%%%%%%%

\section{Conclusion}
\addcontentsline{toc}{section}{Conclusion}

À l’issue de ce premier sprint, les fondations du système sont en place : authentification sécurisée par JWT et bcrypt, connexion via Google OAuth, gestion du profil, et premières interfaces mobiles opérationnelles. Ces éléments sont indispensables pour la suite ; tout ce qui sera développé dans les sprints suivants repose sur une gestion fiable de l’identité utilisateur.

Le sprint 2 aborde désormais la partie la plus technique du projet : la vision par ordinateur et le calcul du score de concentration. 